\begin{abstractEN}
In recent years, deep learning has achieved remarkable progress in video action understanding and analysis tasks within the general domain.
However, a pronounced domain gap persists when migrating such technologies to the healthcare sector.
In clinical practice, compared with natural video analysis, medical video analysis is frequently plagued by low image contrast, subtle motion amplitudes, and highly imbalanced data distributions,
which make it difficult for deep models to learn effective visual representations for accurate action analysis.
Therefore, exploring how to deeply integrate medical domain priors to realize more efficient medical video action understanding is of great significance for the practical deployment of smart healthcare systems.
Focusing on two typical types of medical video data,
namely videofluoroscopic swallowing study (VFSS) videos for auxiliary diagnosis and glaucoma surgery videos for intraoperative navigation,
this paper thoroughly investigates the core challenges of medical video action understanding in complex clinical scenarios as follows:
(1) For temporal micro-action localization in VFSS videos,
the low contrast of X-ray images and the blurring and occlusion of key anatomical structures
easily cause weak motion signals to be overwhelmed by static background noise,
making it hard for models to capture precise temporal boundaries;
(2) For online action recognition in glaucoma surgery videos,
the significant imbalance in the duration of surgical phases and the high visual similarity among different surgical operations
easily lead to missed detections and false judgments during real-time inference.

To address the above challenges, the research contents of this dissertation are presented as follows:

For VFSS video analysis,
this dissertation proposes a novel framework named \sexyname
to tackle the problems of blurred anatomical structures and obscured weak target features in temporal micro-action localization tasks.
This framework introduces a skeleton heatmap modality based on clinical priors to provide the model with robust geometric structural references,
thus guiding it to learn more interference-resistant visual representations.
Furthermore, to ensure efficient calibration and enhancement of multimodal features,
the proposed method employs a Cross-modal Feature Calibration module (\crossblockshort)
and a Channel-level Enhancement mechanism (\channelblockshort).
These modules guide spatial focus through skeleton features and explicitly strip redundant static background signals,
thereby achieving effective amplification of subtle motion signals.
Experimental results on the VFSS benchmark dataset demonstrate that the proposed method
improves the mean Average Precision (mAP) by 5.3\% compared with state-of-the-art methods,
and it exhibits excellent generalization performance on an external validation set.

For glaucoma intraoperative navigation,
this dissertation proposes an online action detection framework named \sexytwo
to address the issues of temporal logic confusion and imbalanced data distribution in online surgical action recognition tasks.
By leveraging the structural priors of surgical procedures through a Dynamic Logic Evolution module (\graphmoduleshort),
this method enhances the model's ability to perceive surgical state transitions and action durations,
thus guiding it to learn more logically coherent scene representations.
Additionally, a Class Boundary Weighted Loss (\lossshort) is proposed to realize balanced learning of features for long and short-duration actions.
Extensive experiments on the self-constructed MIGS dataset and the public OphNet dataset
show that the proposed method achieves an improvement of 0.9\% and 1.9\% in online recognition accuracy,
compared with state-of-the-art methods on the two datasets respectively.

\end{abstractEN}

\keywordsEN{Medical Video Understanding; Videofluoroscopic Swallowing Study; Temporal Micro-action Localization; Glaucoma Surgery Video; Online Surgical Action Detection}